% Standalone section: requires amsmath,amssymb,amsthm.
% The parent document should provide theorem, lemma, corollary, remark environments.
\section{An electrical transport theorem for real zeros}\label{sec:electrical}

Fix a connected undirected graph on $n\geq2$ labelled vertices, with positive
edge conductances.  Parallel edges may be combined by adding conductances.
Let $L$ be its weighted Laplacian, $L^+$ its Moore--Penrose inverse, and
\[
 \tau(G)=\sum_{T\text{ spanning tree}}\prod_{e\in T}w_e,
 \qquad
 R_{ij}=(e_i-e_j)^{\mathsf T}L^+(e_i-e_j).
\]
Thus $R_{ij}$ is the effective resistance between vertices $i$ and $j$.
At vertex $i$, place a monic real-rooted polynomial $p_i$ of degree
$d_i\geq1$, and put $N=d_1+\cdots+d_n$.
Repeated roots, including roots shared by different vertices, are allowed.

A rooted spanning forest consists of a spanning forest together with a
choice of one root in each of its tree components.  Write $\mathcal R(F)$
for its roots.  Define the following polynomial without any divisions:
\begin{equation}\label{eq:forest-polynomial}
 \mathcal P_G(z)=
 \sum_{F\text{ rooted spanning forest}}
 \left(\prod_{e\in F}w_e\right)
 \left(\prod_{i\in\mathcal R(F)}p_i'(z)\right)
 \left(\prod_{i\notin\mathcal R(F)}p_i(z)\right).
\end{equation}
The empty edge set is included, and contributes $\prod_i p_i'$.
The constructions below show that $\mathcal P_G$ is never the zero
polynomial and has degree $N-1$.

List the roots of the input polynomials, with multiplicity and in vertex
blocks, as $r_1,\ldots,r_N$.  Set
\[
 \Lambda=\operatorname{diag}(r_1,\ldots,r_N),\qquad
 V_{ai}=\begin{cases}1&\text{root }a\text{ belongs to }p_i,\\0&\text{otherwise},\end{cases}
 \qquad Q=I_N-\frac{\boldsymbol 1_N\boldsymbol 1_N^{\mathsf T}}{N}.
\]
Here $V$ has $N$ rows and $n$ columns.  Let
$H=\boldsymbol 1_N^\perp$ and use restriction to $H$ whenever a determinant
or eigenvalue of a compressed matrix is taken.  Define
\[
 A_\infty=(Q\Lambda Q)|_H,\qquad
 B_G=(QV L^+V^{\mathsf T}Q)|_H,\qquad K_G=A_\infty-B_G.
\]

\begin{theorem}[Exact compression and electrical transport]\label{thm:electrical}
The following identity holds for every complex $z$:
\begin{equation}\label{eq:compression-identity}
 \boxed{\quad\mathcal P_G(z)=N\tau(G)\det_H(zI_H-K_G).\quad}
\end{equation}
In particular $\mathcal P_G$ has exactly $N-1$ real zeros, with multiplicity.
Let these zeros be $\lambda_1\leq\cdots\leq\lambda_{N-1}$, and let
$\beta_1\leq\cdots\leq\beta_{N-1}$ be the zeros of
$(\prod_i p_i)'$.  Then
\begin{equation}\label{eq:root-upper-bound}
 0\leq\beta_k-\lambda_k\leq\|B_G\|\qquad(1\leq k\leq N-1),
\end{equation}
and
\begin{equation}\label{eq:resistance-sum}
 \boxed{\quad
 \sum_{k=1}^{N-1}(\beta_k-\lambda_k)
 =\frac1N\sum_{i<j}d_i d_jR_{ij}.
 \quad}
\end{equation}
Consequently the empirical probability measures
\[
 \mu_G=\frac1{N-1}\sum_k\delta_{\lambda_k},\qquad
 \nu=\frac1{N-1}\sum_k\delta_{\beta_k}
\]
satisfy the exact optimal-transport formula
\begin{equation}\label{eq:wasserstein-exact}
 \boxed{\quad W_1(\mu_G,\nu)
 =\frac{\sum_{i<j}d_i d_jR_{ij}}{N(N-1)}.\quad}
\end{equation}
Here $W_1$ denotes the minimum average distance needed to transport one
probability measure into the other.  The value in
\eqref{eq:wasserstein-exact} is independent of every input root location.
The bound $W_\infty(\mu_G,\nu)\leq\|B_G\|$ is sharp over the permitted
input polynomials, even with the graph and the degrees fixed.
\end{theorem}

\begin{proof}
We first establish the elementary forest determinant identity
\begin{equation}\label{eq:forest-det}
 \det(L+\operatorname{diag}(x_1,\ldots,x_n))
 =\sum_{F\text{ rooted spanning forest}}
   \left(\prod_{e\in F}w_e\right)\prod_{i\in\mathcal R(F)}x_i.
\end{equation}
One proof adds a new vertex $0$ and an edge of weight $x_i$ from $0$ to
each vertex $i$.  The matrix-tree theorem identifies the left side with
the spanning-tree polynomial of the enlarged graph.  Removing $0$ from
such a tree leaves a spanning forest; each component has exactly one
vertex formerly adjacent to $0$, which becomes its root.  This is a
weight-preserving bijection.  Since it proves a polynomial identity, no
sign assumption on the variables $x_i$ is required.

Away from the input roots, substitute $x_i=p_i'/p_i$ in
\eqref{eq:forest-det} and multiply by $p=\prod_i p_i$.  This gives
\begin{equation}\label{eq:meromorphic-determinant}
 \mathcal P_G(z)=p(z)
 \det\left(L+\operatorname{diag}\left(\frac{p_i'(z)}{p_i(z)}\right)\right).
\end{equation}
By the definition of $V$ and $\Lambda$,
\[
 V^{\mathsf T}(zI_N-\Lambda)^{-1}V
  =\operatorname{diag}\left(\frac{p_i'(z)}{p_i(z)}\right).
\]
The Schur complement of the upper-left block therefore yields the
polynomial identity
\begin{equation}\label{eq:bordered-determinant}
 \mathcal P_G(z)=
 \det\begin{pmatrix}zI_N-\Lambda&V\\-V^{\mathsf T}&L\end{pmatrix}.
\end{equation}
It suffices to establish this first for $z$ outside the finite input
zero set, because both sides are polynomials.

Choose an orthonormal basis of $\mathbb R^n$ starting with
$u=\boldsymbol1_n/\sqrt n$ and otherwise consisting of Laplacian
eigenvectors.  Eliminate the positive-eigenvalue block of $L$ in
\eqref{eq:bordered-determinant}.  The result is
\[
 \det{}'(L)
 \det\begin{pmatrix}
 zI_N-\Lambda+VL^+V^{\mathsf T}&Vu\\
 -u^{\mathsf T}V^{\mathsf T}&0
 \end{pmatrix},
\]
where $\det{}'(L)$ is the product of the $n-1$ positive eigenvalues.
The identity $Vu=\boldsymbol1_N/\sqrt n$ gives
$\|Vu\|^2=N/n$ and $(Vu)^\perp=H$.
For any $N\times N$ matrix $C$ and any nonzero real vector $b$,
\[
 \det\begin{pmatrix}C&b\\-b^{\mathsf T}&0\end{pmatrix}
   =\|b\|^2\det_{b^\perp}(P_{b^\perp}C|_{b^\perp}).
\]
Indeed, in a basis ending with $b/\|b\|$, expansion along the last row
and then the last column gives this equality, including its positive
sign.  Finally $\det{}'(L)=n\tau(G)$ by the matrix-tree theorem.
Together these facts prove \eqref{eq:compression-identity}.

The matrix $K_G$ is real symmetric.  Also $B_G\succeq0$ because
$L^+\succeq0$.  The usual compression identity for derivatives gives
\begin{equation}\label{eq:differentiator}
 \det_H(zI_H-A_\infty)=\frac1N p'(z).
\end{equation}
For completeness, if $e=\boldsymbol1_N/\sqrt N$, the cofactor formula
for a compression says that the left side is
$p(z)e^{\mathsf T}(zI_N-\Lambda)^{-1}e$ whenever the inverse exists.
This is $p(z)N^{-1}\sum_a(z-r_a)^{-1}=p'(z)/N$; polynomial
continuation proves the identity everywhere.  The min--max principle
applied to
$A_\infty-\|B_G\|I_H\preceq K_G\preceq A_\infty$
now proves \eqref{eq:root-upper-bound}.

Put $d=(d_1,\ldots,d_n)^{\mathsf T}$ and
$D=\operatorname{diag}(d_1,\ldots,d_n)$.  Since
$V^{\mathsf T}V=D$ and $V^{\mathsf T}\boldsymbol1_N=d$,
\begin{align*}
 \sum_k(\beta_k-\lambda_k)
 &=\operatorname{tr}B_G
 =\operatorname{tr}\left[L^+\left(D-\frac{dd^{\mathsf T}}N\right)\right]\\
 &=\sum_i d_i(L^+)_{ii}-\frac1N d^{\mathsf T}L^+d
 =\frac1N\sum_{i<j}d_i d_jR_{ij}.
\end{align*}
The last equality follows by expanding
$R_{ij}=(L^+)_{ii}+(L^+)_{jj}-2(L^+)_{ij}$.
For equal-weight measures on the line, an optimal coupling matches
the atoms in increasing order.  Alternatively, the ordered matching
attains the difference of the means here because
$\lambda_k\leq\beta_k$, and every coupling has average distance at
least the absolute difference of the means.  This proves
\eqref{eq:wasserstein-exact} directly.

The ordered matching also proves the $W_\infty$ bound.  To prove
sharpness choose $p_i(z)=(z-c)^{d_i}$ for a common real $c$.
Then $A_\infty=cI_H$; hence the largest displacement is exactly
the largest eigenvalue of $B_G$, namely $\|B_G\|$.
\end{proof}

\begin{corollary}[Electrical barycentre law]\label{cor:barycentre}
The mean of the $N-1$ graph zeros is
\begin{equation}\label{eq:barycentre}
 \frac1{N-1}\sum_k\lambda_k
  =\frac1N\sum_{a=1}^N r_a
   -\frac{\sum_{i<j}d_i d_jR_{ij}}{N(N-1)}.
\end{equation}
\end{corollary}
\begin{proof}
The trace of $A_\infty$ is
$\operatorname{tr}\Lambda-N^{-1}\boldsymbol1_N^{\mathsf T}\Lambda\boldsymbol1_N
=(N-1)N^{-1}\sum_a r_a$.  Substitute this in
\eqref{eq:resistance-sum}.
\end{proof}

\section{Adding one edge: exact root motion}\label{sec:edge-motion}

An edge may be strengthened or inserted.  Fix distinct vertices $a,b$,
put $b_e=e_a-e_b$, and let $L_h=L+h b_e b_e^{\mathsf T}$ for $h\geq0$.
Define quantities using the original connected graph:
\[
 R_e=b_e^{\mathsf T}L^+b_e>0,\qquad
 q_e=QV L^+b_e\in H,\qquad
 \alpha(h)=\frac{h}{1+hR_e}.
\]

\begin{theorem}[Exact edge law]\label{thm:edge-law}
The associated symmetric matrices satisfy
\begin{equation}\label{eq:rank-one-edge}
 \boxed{\quad K_{G_h}=K_G+\frac{h}{1+hR_e}q_e q_e^{\mathsf T}.\quad}
\end{equation}
In particular the ordered graph zeros obey
\[
 \lambda_k(G)\leq\lambda_k(G_h)\leq\lambda_{k+1}(G)
 \quad(1\leq k<N-1),\qquad
 \lambda_{N-1}(G)\leq\lambda_{N-1}(G_h).
\]
Their total displacement is exactly
\begin{equation}\label{eq:edge-total-displacement}
 \boxed{\quad
 \sum_{k=1}^{N-1}\bigl(\lambda_k(G_h)-\lambda_k(G)\bigr)
 =\frac{h}{1+hR_e}\|q_e\|^2.
 \quad}
\end{equation}
This is strictly positive for $h>0$ and does not depend on the input
root locations.  Equivalently,
\[
 W_1(\mu_G,\mu_{G_h})
 =\frac{h\,\|q_e\|^2}{(N-1)(1+hR_e)}.
\]
At a simple eigenvalue, for a normalized real eigenvector $u_k(h)$,
\begin{equation}\label{eq:edge-velocity}
 \frac{d\lambda_k(G_h)}{dh}
 =\frac{|u_k(h)^{\mathsf T}q_e|^2}{(1+hR_e)^2}.
\end{equation}
The total displacement is bounded above by $\|q_e\|^2/R_e$ and
approaches this value as the new conductance tends to infinity.
\end{theorem}

\begin{proof}
On $\boldsymbol1_n^\perp$, the matrix $L$ is positive definite.  The
rank-one inverse identity on that subspace gives
\[
 L_h^+=L^+-\frac{h}{1+hR_e}
                 L^+b_e b_e^{\mathsf T}L^+.
\]
Both sides vanish on $\boldsymbol1_n$, so this is also an identity
on $\mathbb R^n$.  Substitution in the definition of $K_G$ proves
\eqref{eq:rank-one-edge}.  Eigenvalue monotonicity and rank-one
interlacing give the displayed inequalities.  Taking traces proves
\eqref{eq:edge-total-displacement}, and ordered matching proves the
$W_1$ formula.

To check strict positivity, suppose $QV L^+b_e=0$ and write
$x=L^+b_e\in\boldsymbol1_n^\perp$.  Then $Vx$ is constant on all
$N$ coordinates.  As every vertex block is nonempty, all coordinates
of $x$ are equal.  Hence $x=0$, contradicting $Lx=b_e\ne0$.
Thus $q_e\ne0$.  Finally, differentiating the rank-one matrix
identity at a simple eigenvalue yields
\eqref{eq:edge-velocity}, since $\alpha'(h)=(1+hR_e)^{-2}$.
\end{proof}

\begin{remark}[Everything in the displacement budget is electrical]
The vector $q_e$ need not be formed in the $N$-dimensional root space:
\[
 \|q_e\|^2=b_e^{\mathsf T}L^+
      \left(D-\frac{dd^{\mathsf T}}N\right)L^+b_e.
\]
Thus the budget is determined by the graph, the chosen edge, and the
input degrees.  Eigenvectors and the actual zero locations determine
how that fixed budget is distributed among individual zeros.
\end{remark}

\section{A monotone bridge between two derivative constructions}

Scale every conductance by $t>0$ and write $\lambda_k(t)$ for the
zeros of the resulting graph polynomial.  Since $(tL)^+=t^{-1}L^+$,
\begin{equation}\label{eq:uniform-scale}
 K_t=A_\infty-t^{-1}B_G.
\end{equation}
All $N-1$ ordered zeros are nondecreasing with $t$, and they tend to the zeros
of $(\prod_i p_i)'$ as $t\to\infty$.  More precisely,
\[
 0\leq\beta_k-\lambda_k(t)\leq\frac{\|B_G\|}{t},\qquad
 W_1(\mu_{tG},\nu)=
 \frac{\sum_{i<j}d_i d_jR_{ij}}{tN(N-1)}.
\]
There is also a precise small-coupling limit.

\begin{corollary}[Small-coupling limit]\label{cor:small-coupling}
Exactly $n-1$ graph zeros tend to $-\infty$ as $t\downarrow0$.
The other $N-n$ zeros tend, with multiplicity and in increasing order,
to the zeros of $\prod_i p_i'$.  When $N=n$, this latter list is empty.
\end{corollary}
\begin{proof}
On $H$, the matrix $B_G$ has kernel $\ker V^{\mathsf T}$.
Indeed $v^{\mathsf T}B_Gv=0$ is equivalent to
$V^{\mathsf T}v\in\operatorname{span}(\boldsymbol1_n)$; because
$v\perp\boldsymbol1_N$, the sum of the coordinates of
$V^{\mathsf T}v$ is zero, so the latter vector must vanish.
The map $V$ has rank $n$, and hence $B_G$ has rank $n-1$ on $H$.
In the orthogonal decomposition
$H=\operatorname{ran}B_G\oplus\ker B_G$, write
\[
 K_t=\begin{pmatrix}A_{11}-t^{-1}B_1&A_{12}\\
 A_{21}&A_{22}\end{pmatrix},\qquad B_1\succ0.
\]
The first $n-1$ eigenvalues tend to $-\infty$ by min--max.
For $z$ in any fixed bounded set, the inverse of
$zI-A_{11}+t^{-1}B_1$ is $O(t)$, uniformly as $t\downarrow0$.
The complementary Schur determinant therefore tends to
$\det(zI-A_{22})$.  Thus the remaining $N-n$ eigenvalues tend to
the eigenvalues of $A_{22}$, including multiplicity.  This also follows
by applying polynomial continuity to the Schur determinant and using
interlacing to keep these eigenvalues bounded.

Finally $\ker V^{\mathsf T}$ is the direct sum, over the vertex blocks,
of the subspaces whose coordinates sum to zero.  The compression of
$\Lambda$ to these subspaces has characteristic polynomial
$\prod_i(p_i'/d_i)$ by \eqref{eq:differentiator} applied separately
to each block.  These are exactly the asserted finite limiting zeros.
\end{proof}

\section{An exact three-vertex example}\label{sec:exact-example}

Put $p_1(z)=z+1$, $p_2(z)=z$, and $p_3(z)=z-1$ on a three-vertex
path, with each of its two edges having conductance $t>0$.  Enumerating
the rooted forests gives
\[
 \mathcal P_t(z)=1+4tz+t^2(3z^2-1),\qquad
 \lambda_\pm(t)=\frac{-2\pm\sqrt{1+3t^2}}{3t}.
\]
Both zeros increase to $\pm1/\sqrt3$, the zeros of
$(z^3-z)'$.  The three pairwise effective resistances are
$1/t,1/t,2/t$.  The exact transport distance is therefore
\[
 W_1(\mu_t,\nu)=\frac{2}{3t}.
\]
At $t=1$, the graph polynomial is $z(3z+4)$, whose zeros are
$-4/3$ and $0$.  Thus the real-zero theorem does not assert containment
in the interval of the original roots: one new zero lies to the left
of all three input roots.  What is controlled exactly is its ordered
displacement from the zeros of the derivative of their product.

\paragraph{Scope.}
All trace, barycentre and transport formulas in this section concern
finite polynomials.  Applying a formally infinite sum to completed
$L$-functions would require a separate convergence or renormalization
theorem.  The entire-function real-zero and local collision results
do not require such a sum and can be proved directly from logarithmic
derivatives.
